\setcounter{chapter}{-1}
\chapter{Глава для проверки стиля оформления}
\section{Страница для проверки стиля оформления}
\subsection{Непосредственно стиль оформления}
\begin{definition}{Определение}
    \textit{Пиво} -- слабоалкогольный напиток, получаемый спиртовым брожением солодового сусла (чаще всего на основе ячменя) с помощью пивных дрожжей, обычно с добавлением хмеля.
\end{definition}

\begin{statement}
    Существует около тысячи сортов пива.
\end{statement}
\begin{sproof}
    Можете загуглить.
\end{sproof}

\begin{lemma} {Основная лемма о пиве}
    Неверно, что пива не существует
\end{lemma}
\begin{lproof}
    Очевидно.
\end{lproof}

\begin{theorem} {Теорема о пиве}
    Пиво существует
\end{theorem}
\begin{tproof}
    Пусть это не так. Тогда пива не существует, что неверно из основной леммы о пиве. \\
    Таким образом, пиво существует.
\end{tproof}

\begin{note}{}
    Пиво может не существовать, если оно закончилось.
\end{note}

\begin{conseq}
    Если пиво существует, то оно не закончилось.
\end{conseq}